\originalchapter{22.02 作业五:角动量耦合}
题源: Spring 2012 Problem Set 5, 2012年3月26日. 原题所带示范与新增步骤分别说明.
\begin{exercise}\label{ex:2202-5-1}
原题 1 (已解). 两核子自旋均 $1/2$, 轨道 $l=0$. (a) 每核子多少自旋态? (b) 不耦合体系多少态? (c) 耦合表象用什么量子数?
\end{exercise}
\begin{solution}
官方示范重编. (a) 每个两态, $m_s=\pm1/2$. (b) 四态 $\uparrow\uparrow,\uparrow\downarrow,\downarrow\uparrow,\downarrow\downarrow$. (c) $s,m_s,s_1,s_2$, 其中 $s=0,1$, 对应一单重态加三重态. 指定 $s_1=s_2=1/2$ 后可省略这两个固定标签.
\end{solution}
\begin{exercise}\label{ex:2202-5-2}
原题 2. $\mathbf L=\mathbf L_1+\mathbf L_2$. (a) 证明 $[L_1^2,L^2]=0$, 求 $[L_{1z},L^2]$ (原题已解). (b) 证明 $[L_x,L^2]=0$ 并解释 y、z 同样.
\end{exercise}
\begin{solution}
(a) 官方示范重编: $L^2=L_1^2+L_2^2+2\mathbf L_1\cdot\mathbf L_2$. 同粒子 $L_1^2$ 与分量对易, 异粒子算符相互对易, 第一式为零. 第二式为 $2\ii\hbar(L_{1y}L_{2x}-L_{1x}L_{2y})$, 通常非零. (b) 展开 $[L_x,L_y^2]=\ii\hbar(L_zL_y+L_yL_z)$, $[L_x,L_z^2]=-\ii\hbar(L_yL_z+L_zL_y)$, 两项相消, $[L_x,L_x^2]=0$. 循环换指标给另两式.
\end{solution}
\begin{exercise}\label{ex:2202-5-3}
原题 3. 状态 $|l=7,m_x=4\rangle$. (a) 测 $L_x$ 结果? (b) 测 $L_y$ 可能值? (c) 用角动量不确定关系讨论 $\langle L_y\rangle,\langle L_z\rangle$ 是否一致. (d) 假定 $\Delta L_y=\Delta L_z$, 求 $\Delta L_y$.
\end{exercise}
\begin{solution}
(a) $4\hbar$ 确定. (b) $m_y\hbar$, $m_y=-7,\ldots,7$, 具体非零概率由转动系数决定. (c) 对绕 x 的转动, 该态只获整体相位, 横向均值必须为零, 因而两者均零. 不确定关系正确方向为 $\Delta L_a\Delta L_b\ge\hbar|\langle L_c\rangle|/2$, 原题把 $\ge$ 印成 $\le$. $\Delta L_x=0$ 与横向均值零相容. (d) $\langle L^2\rangle=56\hbar^2$, $\langle L_x^2\rangle=16\hbar^2$, 余两项相等各 $20\hbar^2$, 所以 $\Delta L_y=\sqrt{20}\hbar$.
\end{solution}
\begin{exercise}\label{ex:2202-5-4}
原题 4. 对 $|l,m_z\rangle$, (a) 求 $L_\pm$ 的作用及均值; (b) 用此证明 $\langle L_x\rangle=0$ 对所有 $L_z,L^2$ 共同本征态成立.
\end{exercise}
\begin{solution}
作用为第 \ref{ch:angular} 章梯算符公式, 不同 m 的本征态正交, 所以 $\langle L_\pm\rangle=0$. $L_x=(L_++L_-)/2$, $L_y=(L_+-L_-)/(2\ii)$, 于是两个均值都零. 最高最低 m 的梯作用自身为零, 结论仍成立.
\end{solution}
\begin{exercise}\label{ex:2202-5-5}
原题 5. 氢原子电子态 $R_{21}(r)[Y_{10}|\downarrow\rangle/\sqrt3+\sqrt{2/3}Y_{1,-1}|\uparrow\rangle]$. (a) 求可能 j (已解). (b) 对照耦合表示 $R_{21}[2\sqrt2|3/2,-1/2\rangle/3-|1/2,-1/2\rangle/3]$, 哪个是耦合表象, 如何求转换, 是否与 (a) 一致? (c) 给 $L^2,L_z,S_z,J^2,J_z$ 的结果和概率. (d) 两 p 电子耦合态 $|2,-1;1,1\rangle$, 可能 $m_{1z},m_{2z}$? (已解) (e) 对 $|2,-2;1,1\rangle$, 测 $L_{1z}$ 的值、概率以及二者都为 $-\hbar$ 的联合概率.
\end{exercise}
\begin{solution}
(a) 官方示范: $j=1/2,3/2$. (b) 第二式按 j、$m_j$ 标记, 是耦合表象. $|3/2,-1/2\rangle=\sqrt{2/3}|m_l=0,\downarrow\rangle+\sqrt{1/3}|m_l=-1,\uparrow\rangle$, 可选 $|1/2,-1/2\rangle=\sqrt{1/3}|0,\downarrow\rangle-\sqrt{2/3}|-1,\uparrow\rangle$. 求内积给系数 $2\sqrt2/3,-1/3$, 平方和为1, 只含允许的两个 j.

(c) $L^2=2\hbar^2$ 确定; $L_z=0,-\hbar$ 的概率为 $1/3,2/3$; $S_z=-\hbar/2,+\hbar/2$ 的概率为 $1/3,2/3$; $J^2=15\hbar^2/4,3\hbar^2/4$ 的概率为 $8/9,1/9$; $J_z=-\hbar/2$ 确定. (d) 两种 $(m_{1z},m_{2z})=(-1,0),(0,-1)$, 对称 $j=2$ 态各概率 $1/2$. (e) 只有 $(-1,-1)$, 所以单次 $L_{1z}=-\hbar$ 与联合结果都概率1. 这只是轨道部分的指定态, 若把两粒子视为同一原子中的相同电子, 完整空间--自旋态还应满足交换反对称性.
\end{solution}
